\documentclass[11pt,a4paper]{article}
\usepackage[T1]{fontenc}
\usepackage[utf8]{inputenc}
\usepackage[english]{babel}
\usepackage{amsmath,amssymb}
\usepackage[margin=2cm]{geometry}
\setlength{\parindent}{0pt}
\setlength{\emergencystretch}{3em}
\allowdisplaybreaks

\title{About \texttt{index\_notation\_beyond.tex}}
\author{}
\date{}
\begin{document}
\maketitle
This note explains what the test data in \texttt{index\_notation\_beyond.tex} is used for. That file is input for the engine in \texttt{main.c}; it is read by the test \texttt{test\_index\_identities.c} in this folder.

\begin{enumerate}
\item \textbf{Purpose.} The file holds 84 statements of tensor calculus, with a Kronecker delta, a Levi-Civita symbol or a declared metric, that are outside the language the engine reads today (identities with indices in a fixed dimension, with or without a declared diagonal metric, with tensors given by their components, with Greek indices and with values of the indices that start from 0 or from 1). The engine must answer \emph{unsupported}, with a reason, for each of them: it may neither prove nor refute a statement it cannot read in the sense its author meant.

\item \textbf{What is in the file.} Statements without a declared dimension ($\delta_{ii} = 3$ alone holds only in three dimensions); statements that break the summation convention (an index three times in a product, different free indices in the two sides, a number other than $0$ where a tensor with a free index is expected); symbols with the wrong number of indices; notation that needs a metric that the statement does not declare (the same tensor, or $\varepsilon$, with upper and lower indices, as in $g_{ij} g^{jk} = \delta_{i}^{k}$ and $\varepsilon^{ijkl}\varepsilon_{ijkl} = -24$, ); statements that declare a metric but leave out the value of $\varepsilon$, give a metric that is not diagonal, or whose entries are identically zero, have indices or use a trigonometric function, or contract two lower indices; hypotheses that give a tensor by its components without writing its indices, with a table that is not square or does not have the size of the metric, with an element that names a given tensor, a symbol or an index letter, or with two numbers that no sign separates; Greek indices whose values are not declared, the index $0$ without a declaration, two declarations that do not agree, $\delta$ used as an index; derivatives and Christoffel symbols; vectors, powers, relations other than equality, chains of equalities. Several are true in their intended reading, for example $\varepsilon_{\mu\nu\rho\sigma}\varepsilon^{\mu\nu\rho\sigma} = -24$ in Minkowski space: read without the metric that its author had in mind it would be declared false.

\item \textbf{What the engine needs in order to pass.} To refuse each statement while reading it. When a later stage makes one of them readable, it must be moved to the file of the statements that have to be proved, together with the conventions it declares.

\item \textbf{Format expected by the test.} One \texttt{\textbackslash item} per statement, with the formula written as \texttt{\$\textbackslash displaystyle ...\$}.

\item \textbf{How to run it.} From the top folder of the repository:
\begin{verbatim}
cc -std=c11 -O2 -Wall -Wextra -Werror -o /tmp/test_index_identities \
    tests/tensor_calculus/test_index_identities.c && /tmp/test_index_identities
\end{verbatim}

\item \textbf{Expected result.} The line \texttt{PASS: 84 of 84 statements outside the language are unsupported, each with its reason}.
\end{enumerate}
\end{document}
